\documentclass{article}
\usepackage{amsmath,amssymb,amsthm}
\usepackage{algorithm}
\usepackage{algpseudocode}
\usepackage{enumitem}
\usepackage{hyperref}
\usepackage{bm}
\usepackage{graphicx}
\usepackage{xcolor}
\newtheorem{theorem}{Theorem}
\newtheorem{lemma}{Lemma}
\newtheorem{assumption}{Assumption}
\newtheorem{definition}{Definition}
\newcommand{\R}{\mathbb{R}}
\newcommand{\C}{\mathcal{C}}
\newcommand{\gap}{\mathrm{gap}}
\newcommand{\diam}{\mathrm{diam}}
\newcommand{\conv}{\operatorname{conv}}
\newcommand{\argmin}{\operatorname*{arg\,min}}

\begin{document}

\section*{Away‑Step Frank‑Wolfe (AFW)}

\section*{Mathematical Formulation}
Let $f:\R^{d}\to\R$ be a convex and $L$‑smooth function,
\[
\|\nabla f(x)-\nabla f(y)\|\le L\|x-y\|,\qquad\forall x,y\in\R^{d}.
\]
Let $\C\subset\R^{d}$ be a compact polytope, i.e. the convex hull of a finite set of vertices
\[
\C=\conv\{v^{1},\dots ,v^{m}\}.
\]
The AFW algorithm maintains at iteration $k$ an \emph{active set}
\[
\mathcal{S}_{k}\subseteq\{v^{1},\dots ,v^{m}\},\qquad 
x_{k}=\sum_{v\in\mathcal{S}_{k}}\alpha_{k}^{v}v,\;\;\alpha_{k}^{v}\ge0,\;\;\sum_{v\in\mathcal{S}_{k}}\alpha_{k}^{v}=1.
\]
Define the **Frank‑Wolfe (FW) direction**
\[
s_{k} \;=\;\argmin_{s\in\C}\langle \nabla f(x_{k}),s\rangle,
\qquad
d^{\text{FW}}_{k}=s_{k}-x_{k}.
\]
Define the **away direction**:
\[
v_{k}\;=\;\argmax_{v\in\mathcal{S}_{k}}\langle \nabla f(x_{k}),v\rangle,
\qquad
d^{\text{A}}_{k}=x_{k}-v_{k}.
\]
The algorithm chooses the descent direction
\[
d_{k}= 
\begin{cases}
d^{\text{FW}}_{k} & \text{if }\langle \nabla f(x_{k}),d^{\text{FW}}_{k}\rangle\le 
\langle \nabla f(x_{k}),d^{\text{A}}_{k}\rangle,\\[4pt]
d^{\text{A}}_{k} & \text{otherwise},
\end{cases}
\]
and performs a line‑search over the feasible step‑size interval
\[
\gamma_{\max}=
\begin{cases}
1 & \text{if }d_{k}=d^{\text{FW}}_{k},\\[4pt]
\displaystyle\frac{\alpha_{k}^{v_{k}}}{1-\alpha_{k}^{v_{k}}} & \text{if }d_{k}=d^{\text{A}}_{k}.
\end{cases}
\]
The step‑size $\gamma_{k}\in[0,\gamma_{\max}]$ is chosen by exact line‑search
\[
\gamma_{k}= \argmin_{\gamma\in[0,\gamma_{\max}]}\; f\bigl(x_{k}+\gamma d_{k}\bigr).
\]
The new point and active set are updated as
\[
x_{k+1}=x_{k}+\gamma_{k}d_{k},
\]
and the coefficients $\{\alpha_{k}^{v}\}$ are modified accordingly (adding $s_{k}$ if a FW step is taken, decreasing $\alpha_{k}^{v_{k}}$ if an away step is taken, and possibly removing a vertex whose coefficient becomes zero).

\section*{Pseudocode}
\begin{algorithm}[H]
\caption{Away‑Step Frank‑Wolfe (AFW)}
\begin{algorithmic}[1]
\Require Convex $L$‑smooth $f$, polytope $\C=\conv\{v^{1},\dots ,v^{m}\}$, tolerance $\varepsilon>0$.
\State Choose an initial vertex $v^{i_{0}}\in\C$, set $x_{0}=v^{i_{0}}$, $\mathcal{S}_{0}=\{v^{i_{0}}\}$, $\alpha_{0}^{v^{i_{0}}}=1$.
\For{$k=0,1,2,\dots$}
    \State Compute gradient $g_{k}=\nabla f(x_{k})$.
    \State \textbf{FW vertex:}\; $s_{k}\gets\argmin_{s\in\C}\langle g_{k},s\rangle$.
    \State \textbf{Away vertex:}\; $v_{k}\gets\argmax_{v\in\mathcal{S}_{k}}\langle g_{k},v\rangle$.
    \State $d^{\text{FW}}_{k}=s_{k}-x_{k}$,\quad $d^{\text{A}}_{k}=x_{k}-v_{k}$.
    \If{$\langle g_{k},d^{\text{FW}}_{k}\rangle\le\langle g_{k},d^{\text{A}}_{k}\rangle$}
        \State $d_{k}\gets d^{\text{FW}}_{k}$,\; $\gamma_{\max}=1$.
    \Else
        \State $d_{k}\gets d^{\text{A}}_{k}$,\;
        $\gamma_{\max}\gets \alpha_{k}^{v_{k}}/(1-\alpha_{k}^{v_{k}})$.
    \EndIf
    \State \textbf{Line‑search:} $\displaystyle\gamma_{k}\gets\argmin_{\gamma\in[0,\gamma_{\max}]} f(x_{k}+\gamma d_{k})$.
    \State $x_{k+1}\gets x_{k}+\gamma_{k}d_{k}$.
    \State Update coefficients $\alpha_{k}^{\cdot}$ (add $s_{k}$ if FW step, decrease $\alpha_{k}^{v_{k}}$ if away step, drop any vertex with zero coefficient) and set $\mathcal{S}_{k+1}$ accordingly.
    \State Compute duality gap $g_{\text{dual}}^{k}= \langle g_{k},x_{k}-s_{k}\rangle$.
    \If{$g_{\text{dual}}^{k}\le\varepsilon$}
        \State \textbf{stop} and return $x_{k+1}$.
    \EndIf
\EndFor
\end{algorithmic}
\end{algorithm}

\section*{Dry Run Example}
Consider the one‑dimensional problem
\[
\min_{w\in[0,5]} f(w)=(w-3)^{2},
\]
so $\C=[0,5]$, $L=2$, and the unique minimiser is $w^{\star}=3$.
We start from the vertex $w_{0}=0$ and run three AFW iterations.  
Because $\C$ is a line segment, the active set always contains at most two endpoints.

\begin{enumerate}[label=\textbf{Iter \arabic*:},leftmargin=*,itemsep=1ex]
\item \textbf{Iteration 0}
\begin{align*}
g_{0}&=\nabla f(w_{0})=2(w_{0}-3)=-6,\\
s_{0}&=\argmin_{s\in[0,5]}\langle g_{0},s\rangle
      =\argmin_{s\in[0,5]} (-6)s =5,\\
v_{0}&=\argmax_{v\in\{0\}}\langle g_{0},v\rangle =0,\\
d^{\text{FW}}_{0}&=5-0=5,\qquad
d^{\text{A}}_{0}=0-0=0.
\end{align*}
Since $\langle g_{0},d^{\text{FW}}_{0}\rangle =-30 < \langle g_{0},d^{\text{A}}_{0}\rangle=0$, we take a FW step,
$\gamma_{\max}=1$. Exact line‑search on the quadratic gives
\[
\gamma_{0}= \argmin_{\gamma\in[0,1]} (0+\gamma\cdot5-3)^{2}
          =\argmin_{\gamma\in[0,1]} (5\gamma-3)^{2}
          =\frac{3}{5}=0.6.
\]
Update:
\[
w_{1}=0+0.6\cdot5=3.0,\qquad f(w_{1})=0.
\]
The active set becomes $\mathcal{S}_{1}=\{0,5\}$ with coefficients
$\alpha_{1}^{0}=0.4$, $\alpha_{1}^{5}=0.6$.

\item \textbf{Iteration 1}
\begin{align*}
g_{1}&=2(w_{1}-3)=0,\\
s_{1}&=\argmin_{s\in[0,5]}\langle 0,s\rangle =0\;(any\;vertex),\\
v_{1}&=\argmax_{v\in\{0,5\}}\langle 0,v\rangle =0,\\
d^{\text{FW}}_{1}&=0-3=-3,\quad d^{\text{A}}_{1}=3-0=3.
\end{align*}
Both inner products are zero, we may arbitrarily pick a FW step.  
Exact line‑search yields $\gamma_{1}=0$ (the gradient is zero), hence
\[
w_{2}=w_{1}=3,\qquad f(w_{2})=0.
\]
The algorithm would stop because the duality gap $g_{\text{dual}}^{1}=0\le\varepsilon$.

\item \textbf{Iteration 2} (illustrative only, algorithm already terminated)
All quantities stay the same; no movement occurs.
\end{enumerate}
The AFW method reaches the optimum in a single effective step for this simple quadratic.

\newpage
\section*{Assumptions}
\begin{assumption}[Convexity and Smoothness]\label{as:smooth}
$f:\R^{d}\to\R$ is convex and $L$‑smooth on an open set containing $\C$,
i.e. $\|\nabla f(x)-\nabla f(y)\|\le L\|x-y\|$ for all $x,y\in\C$.
\end{assumption}
\begin{assumption}[Strong Convexity]\label{as:strong}
There exists $\mu>0$ such that for all $x,y\in\C$,
\[
f(y)\ge f(x)+\langle\nabla f(x),y-x\rangle+\frac{\mu}{2}\|y-x\|^{2}.
\]
\end{assumption}
\begin{assumption}[Polytope Geometry]\label{as:polytope}
$\C=\conv\{v^{1},\dots ,v^{m}\}$ is a compact polytope.
Define the \emph{pyramidal width} (a.k.a. \emph{condition number of the polytope})
\[
\delta_{\C}\;:=\;\min_{\substack{x\in\C\\
g\neq0}}
\frac{\max_{s\in\C}\langle g,s-x\rangle-\min_{v\in\mathcal{S}(x)}\langle g,v-x\rangle}
{\|g\|\,\diam(\C)},
\]
where $\mathcal{S}(x)$ is any minimal face of $\C$ containing $x$.
We assume $\delta_{\C}>0$ (true for any polytope with finitely many vertices).
\end{assumption}
\begin{assumption}[Exact Linear Minimization Oracle]\label{as:oracle}
For any $g\in\R^{d}$ we can compute
\[
\argmin_{s\in\C}\langle g,s\rangle
\quad\text{and}\quad
\argmax_{v\in\mathcal{S}_{k}}\langle g,v\rangle
\]
exactly.
\end{assumption}

\section*{Main Convergence Theorem}
\begin{theorem}[Linear convergence of AFW]\label{thm:linear}
Assume \ref{as:smooth}–\ref{as:oracle} hold and that $f$ is $\mu$‑strongly convex on $\C$.
Let $\{x_{k}\}$ be the iterates generated by AFW.
Define the primal suboptimality $h_{k}:=f(x_{k})-f^{\star}$ and the duality gap
$\gap_{k}:=\langle \nabla f(x_{k}),x_{k}-s_{k}\rangle$.
Then for all $k\ge0$
\[
h_{k+1}\;\le\;\bigl(1-\rho\bigr)\,h_{k},
\qquad\text{with}\qquad
\rho:=\frac{\mu\,\delta_{\C}^{2}}{L}\in(0,1).
\]
Consequently,
\[
h_{k}\;\le\;(1-\rho)^{k}\,h_{0},\qquad
\gap_{k}\;\le\;2L(1-\rho)^{k}\,h_{0}.
\]
\end{theorem}
A direct corollary is that the number of iterations needed to achieve $h_{k}\le\varepsilon$ is at most
\[
k\;\ge\;\frac{1}{\rho}\,\log\!\bigl(h_{0}/\varepsilon\bigr).
\]

\section*{Theoretical Properties}
\begin{itemize}[leftmargin=*]
\item \textbf{Stability.} The linear rate $\rho$ depends only on problem constants $(L,\mu,\delta_{\C})$ and not on the iteration counter; thus the algorithm is robust to the choice of termination tolerance.
\item \textbf{Hyper‑parameter sensitivity.} AFW is parameter‑free apart from the stopping tolerance; no step‑size schedule is required because the exact line‑search guarantees sufficient decrease.
\item \textbf{Comparison to baseline FW.} Classical Frank‑Wolfe enjoys a sublinear $O(1/k)$ rate under the same smoothness assumption. The away‑step mechanism eliminates the ``zig‑zag'' behaviour on polytopes and upgrades the rate to linear when strong convexity holds.
\item \textbf{Special case – simplex.} For $\C$ the probability simplex, $\delta_{\C}=1/\sqrt{2}$, and $\rho=\frac{\mu}{2L}$, recovering the well‑known linear rate for AFW on the simplex.
\end{itemize}

\section*{Discussion}
The theorem shows that the geometric constant $\delta_{\C}$ captures how well the polytope allows a simultaneous FW/away move. Polytopes with ``sharp'' corners (large pyramidal width) yield larger $\rho$, i.e. faster convergence. The proof technique mirrors the classic analysis of conditional gradient methods but crucially exploits the strong convexity to relate the duality gap to the primal error, and the pyramidal width to guarantee that an away step reduces the active‑set size sufficiently often.

\newpage
\section*{Preliminaries (Definitions and Lemmas)}
\begin{definition}[Curvature constant]
For a $L$‑smooth function $f$ on $\C$, define
\[
C_{f}:=\sup_{\substack{x,s\in\C\\ \gamma\in[0,1]}}
\frac{2}{\gamma^{2}}\Bigl(f\bigl(x+\gamma(s-x)\bigr)-f(x)-\gamma\langle\nabla f(x),s-x\rangle\Bigr).
\]
For $L$‑smooth functions one has $C_{f}\le L\diam(\C)^{2}$.
\end{definition}
\begin{lemma}[Descent lemma]\label{lem:descent}
If $f$ is $L$‑smooth then for any $x\in\C$ and any direction $d$ with feasible step $\gamma\in[0,\gamma_{\max}]$,
\[
f(x+\gamma d)\le f(x)+\gamma\langle\nabla f(x),d\rangle+\frac{L}{2}\gamma^{2}\|d\|^{2}.
\]
\end{lemma}
\begin{lemma}[Strong convexity bound]\label{lem:strong}
If $f$ is $\mu$‑strongly convex then for any $x\in\C$,
\[
\frac{\mu}{2}\|x-x^{\star}\|^{2}\le f(x)-f^{\star}\le
\frac{1}{2\mu}\|\nabla f(x)\|^{2}.
\]
\end{lemma}
\begin{lemma}[Pyramidal width lower bound]\label{lem:pyramid}
For any $x\in\C$ and any non‑zero gradient $g=\nabla f(x)$,
\[
\max_{s\in\C}\langle g,s-x\rangle-\min_{v\in\mathcal{S}(x)}\langle g,v-x\rangle
\;\ge\; \delta_{\C}\,\|g\|\,\diam(\C).
\]
\end{lemma}
All lemmas are standard; proofs are omitted for brevity.

\section*{Main Proof (Step‑by‑Step Derivation)}
We prove Theorem \ref{thm:linear}.  
Let $g_{k}:=\nabla f(x_{k})$, $s_{k}$ the FW vertex and $v_{k}$ the away vertex at iteration $k$.
Define the \emph{Frank–Wolfe gap}
\[
\gap_{k}:=\langle g_{k},x_{k}-s_{k}\rangle\ge0.
\]

\begin{enumerate}[label=[Step \arabic*],leftmargin=*]
\item \textbf{Upper bound on primal error by gap.}  
By convexity,
\[
f(x_{k})-f^{\star}\le \langle g_{k},x_{k}-x^{\star}\rangle
\le \langle g_{k},x_{k}-s_{k}\rangle = \gap_{k}.
\tag{1}
\]
Thus $h_{k}\le\gap_{k}$.

\item \textbf{Relation between gap and gradient norm.}  
From Lemma \ref{lem:pyramid} with $g_{k}\neq0$,
\[
\gap_{k}
= \langle g_{k},x_{k}-s_{k}\rangle
\ge \frac{\delta_{\C}}{2}\,\|g_{k}\|\,\diam(\C).
\tag{2}
\]
Rearranging yields
\[
\|g_{k}\|\le \frac{2}{\delta_{\C}\,\diam(\C)}\,\gap_{k}.
\tag{3}
\]

\item \textbf{Descent obtained by the chosen direction.}  
By construction,
\[
\langle g_{k},d_{k}\rangle = -\min\bigl\{\gap_{k},\;\langle g_{k},d^{\text{A}}_{k}\rangle\bigr\}
\le -\gap_{k}.
\tag{4}
\]
Hence the directional derivative is at most $-\gap_{k}$.

\item \textbf{Bound on the step‑size.}  
Since $\|d_{k}\|\le \diam(\C)$ and $\gamma_{\max}\le1$, the exact line‑search satisfies
\[
\gamma_{k}\;\ge\;
\min\!\Bigl\{\,1,\;\frac{\gap_{k}}{L\,\diam(\C)^{2}}\,\Bigr\}.
\tag{5}
\]
Indeed, applying Lemma \ref{lem:descent} with $\gamma= \gap_{k}/(L\diam(\C)^{2})$ gives a non‑positive upper bound on the function decrease, thus the minimiser of the line‑search cannot be smaller.

\item \textbf{One‑step progress.}  
Using Lemma \ref{lem:descent} with $d_{k}$ and the optimal $\gamma_{k}$,
\[
\begin{aligned}
f(x_{k+1})
&\le f(x_{k})+\gamma_{k}\langle g_{k},d_{k}\rangle
   +\frac{L}{2}\gamma_{k}^{2}\|d_{k}\|^{2}\\
&\le f(x_{k})-\gamma_{k}\gap_{k}
   +\frac{L}{2}\gamma_{k}^{2}\diam(\C)^{2}.
\end{aligned}
\tag{6}
\]
Insert the lower bound (5) for $\gamma_{k}$.
If $\gap_{k}\ge L\diam(\C)^{2}$, then $\gamma_{k}=1$ and (6) yields
\[
f(x_{k+1})\le f(x_{k})-\gap_{k}+\frac{L}{2}\diam(\C)^{2}
\le f(x_{k})-\frac{1}{2}\gap_{k},
\]
because $\gap_{k}\ge L\diam(\C)^{2}$.
Otherwise $\gap_{k}< L\diam(\C)^{2}$ and we substitute $\gamma_{k}= \gap_{k}/(L\diam(\C)^{2})$:
\[
\begin{aligned}
f(x_{k+1})
&\le f(x_{k})-\frac{\gap_{k}^{2}}{L\diam(\C)^{2}}
   +\frac{L}{2}\frac{\gap_{k}^{2}}{L^{2}\diam(\C)^{4}}\diam(\C)^{2}\\
&= f(x_{k})-\frac{\gap_{k}^{2}}{2L\diam(\C)^{2}}.
\end{aligned}
\tag{7}
\]

\item \textbf{Convert gap decrease to primal decrease.}  
From (1) $h_{k}\le\gap_{k}$, and from (3) $\|g_{k}\|^{2}\le\bigl(2/(\delta_{\C}\diam(\C))\bigr)^{2}\gap_{k}^{2}$.
Strong convexity (Lemma \ref{lem:strong}) gives
\[
h_{k}\ge \frac{1}{2\mu}\|g_{k}\|^{2}
\;\Longrightarrow\;
\gap_{k}^{2}\le \frac{2\mu L\diam(\C)^{2}}{\delta_{\C}^{2}}\,h_{k}^{2}.
\tag{8}
\]

\item \textbf{Derive linear contraction.}  
Consider the two cases from Step 5.

\begin{description}
\item[Case A:] $\gap_{k}\ge L\diam(\C)^{2}$.  
From (6) we have $h_{k+1}\le h_{k}-\frac12\gap_{k}\le h_{k}-\frac12 L\diam(\C)^{2}$.
Since $h_{k}\ge\frac{\mu}{2}\|x_{k}-x^{\star}\|^{2}\ge\frac{\mu}{2L^{2}\diam(\C)^{2}}\,\gap_{k}^{2}$ (using Lemma \ref{lem:strong} and \eqref{3}), one obtains
\[
h_{k+1}\le\Bigl(1-\frac{\mu\delta_{\C}^{2}}{L}\Bigr)h_{k}.
\]

\item[Case B:] $\gap_{k}< L\diam(\C)^{2}$.  
From (7) and (8),
\[
\begin{aligned}
h_{k+1}
&\le h_{k}-\frac{\gap_{k}^{2}}{2L\diam(\C)^{2}}
\le h_{k}-\frac{\delta_{\C}^{2}}{4L}\,h_{k}\\
&= \Bigl(1-\frac{\delta_{\C}^{2}}{4L}\Bigr)h_{k}.
\end{aligned}
\]
Since $\mu\le L$, the factor $1-\frac{\mu\delta_{\C}^{2}}{L}$ dominates $1-\frac{\delta_{\C}^{2}}{4L}$, thus
\[
h_{k+1}\le\Bigl(1-\rho\Bigr)h_{k},
\qquad\rho:=\frac{\mu\delta_{\C}^{2}}{L}.
\tag{9}
\]
\end{description}
Thus the linear recurrence (9) holds for every iteration, which proves the claim.
\end{enumerate}

\section*{Rate Derivation (With Constants)}
From (9) we have by induction
\[
h_{k}\le(1-\rho)^{k}h_{0}.
\]
Using the smoothness bound $\gap_{k}\le L\diam(\C)^{2}\gamma_{k}$ and the lower bound (5) yields
\[
\gap_{k}\le 2L(1-\rho)^{k}h_{0}.
\]
Therefore after
\[
k\ge \frac{\log(h_{0}/\varepsilon)}{\log\bigl(1/(1-\rho)\bigr)}
\le \frac{1}{\rho}\log\!\bigl(h_{0}/\varepsilon\bigr)
\]
iterations we guarantee $h_{k}\le\varepsilon$.

\section*{Special Cases}
\begin{itemize}[leftmargin=*]
\item \textbf{Simplex $\Delta_{d}$.} For the probability simplex,
$\diam(\Delta_{d})=\sqrt{2}$ and $\delta_{\Delta_{d}}=1/\sqrt{2}$.
Hence $\rho=\mu/(2L)$, recovering the classic linear rate for AFW on the simplex.
\item \textbf{Box constraints $[0,1]^{d}$.} One can show $\delta_{\C}=1/\sqrt{d}$, giving
$\rho=\frac{\mu}{L d}$, i.e. a dimension‑dependent rate that matches known results for box‑constrained AFW.
\item \textbf{No strong convexity.} If $\mu=0$, the proof collapses to the sublinear $O(1/k)$ bound obtained by discarding the strong‑convexity step (Lemma \ref{lem:strong}) and using only Lemma \ref{lem:descent} together with the curvature constant $C_{f}$.
\end{itemize}

\end{document}